\clearpage\thispagestyle{empty}
\pdfbookmark[0]{Quellen zum Bericht über die Sternwarte Seeberg}{seeberg1808-sources}
\begin{center}\textbf{Quellen zum Bericht über die Sternwarte Seeberg}\end{center}
\noindent Diese KI-Übertragung ins heutige Deutsch folgt dem Bericht von Gauß auf S.~498--499 der \href{https://gdz.sub.uni-goettingen.de/id/PPN236007467}{\textit{Werke}, Bd.~VI} (Göttingen, 1874). Gauß stützt sich auf einen Bericht von Bernhard von Lindenau. Der Text erschien ohne Namensnennung unter der Überschrift \emph{Göttingen} in den \textit{Göttingischen gelehrten Anzeigen}, Stück~106, 2. Juli 1808, S.~\href{https://api.digitale-sammlungen.de/iiif/image/v2/bsb10538674_00017/full/1800,/0/default.jpg}{1049}--\href{https://api.digitale-sammlungen.de/iiif/image/v2/bsb10538674_00020/full/1800,/0/default.jpg}{1052}. Das verlinkte Exemplar stammt aus der Bayerischen Staatsbibliothek.
\par\smallskip\noindent Der Erstdruck zeigt bei der Sonnenbeobachtung vom 9. Mai auf \href{https://api.digitale-sammlungen.de/iiif/image/v2/bsb10538674_00018/full/1800,/0/default.jpg}{S.~1050} ein Pluszeichen. In der Werkausgabe steht ein waagerechter Strich mit einem abgesetzten Strich darunter. Die übertragene Tabelle gibt die sichtbare Form aus der Werkausgabe wieder. Das Zeichen bleibt mehrdeutig. KI-Übertragung und KI-Satz.
\clearpage
